% Standalone polished architecture template. Compile this file with pdfLaTeX.
% No external input files, images, bibliography database, or shell escape needed.
% Adapted from Rich Hilliard's templates, copyright 2012--2014, CC BY 3.0.
% This 2026 working edition changes layout, wording, fields, and edition guidance.
% Preserve attribution and identify further changes when reusing this adaptation.
\documentclass[11pt,letterpaper]{article}
% Polished architecture-template style, working edition 1.0, 2026-09-06.
% Template content is adapted from Rich Hilliard's CC BY 3.0 templates.
\RequirePackage[T1]{fontenc}
\RequirePackage[utf8]{inputenc}
\RequirePackage{mathpazo}
\RequirePackage[scaled=0.92]{helvet}
\RequirePackage{microtype}
\RequirePackage[letterpaper,left=0.77in,right=0.77in,top=0.73in,bottom=0.73in,headheight=21pt,headsep=18pt,footskip=25pt]{geometry}
\RequirePackage[table]{xcolor}
\RequirePackage{array,tabularx,booktabs}
\RequirePackage{enumitem}
\RequirePackage[skins,breakable]{tcolorbox}
\RequirePackage{titlesec,fancyhdr,lastpage,needspace}
\RequirePackage{xurl}
\RequirePackage{hyperref}
\definecolor{Navy}{HTML}{193443}
\definecolor{Teal}{HTML}{137980}
\definecolor{Ink}{HTML}{293D47}
\definecolor{Muted}{HTML}{62737D}
\definecolor{Line}{HTML}{D6E2E6}
\definecolor{Pale}{HTML}{F1F6F7}
\definecolor{Tint}{HTML}{EBF5F3}
\color{Ink}
\hypersetup{colorlinks=true,linkcolor=Teal,urlcolor=Teal,citecolor=Teal}
\urlstyle{same}
\setlength{\parindent}{0pt}
\setlength{\parskip}{6pt}
\setlength{\emergencystretch}{1.5em}
\linespread{1.035}
\raggedbottom
\renewcommand{\arraystretch}{1.2}
\setlength{\tabcolsep}{5pt}
\newcolumntype{L}[1]{>{\raggedright\arraybackslash}p{#1}}
\newcolumntype{Y}{>{\raggedright\arraybackslash}X}
\setlist[itemize]{leftmargin=1.25em,itemsep=3pt,topsep=4pt,parsep=0pt}
\setlist[enumerate]{leftmargin=1.55em,itemsep=4pt,topsep=4pt,parsep=0pt}
\titleformat{\section}{\sffamily\bfseries\Large\color{Navy}}{\textcolor{Teal}{\thesection}}{0.65em}{}
\titleformat{\subsection}{\sffamily\bfseries\normalsize\color{Navy}}{}{0pt}{}
\titlespacing*{\section}{0pt}{0pt}{9pt}
\titlespacing*{\subsection}{0pt}{11pt}{4pt}
\providecommand{\TemplateCode}{ARCHITECTURE TEMPLATE}
\providecommand{\TemplateShortTitle}{WORKING EDITION}
\newcommand{\TemplatePageBreak}{\clearpage}
\newcommand{\FillIn}[1]{\textcolor{Teal}{\sffamily\itshape[#1]}}
\newcommand{\RecordID}[1]{\textsf{\textbf{#1}}}
\newcommand{\TableHead}[1]{\textcolor{Navy}{\sffamily\bfseries#1}}
\newcommand{\Lead}[1]{\textcolor{Muted}{#1}\par\vspace{3pt}}
\newcommand{\Repeat}[1]{\textcolor{Teal}{\sffamily\small\textbf{REPEAT}\enspace #1}\par}
\newcommand{\Guide}[1]{\par{\small\textcolor{Muted}{\textbf{Authoring guidance.} #1}}\par}
\newtcolorbox{TemplateNote}[1]{enhanced,colback=Tint,colframe=Tint,
  coltitle=Navy,title={#1},fonttitle=\sffamily\bfseries,boxrule=0pt,arc=1pt,
  borderline west={2pt}{0pt}{Teal},left=10pt,right=10pt,top=6pt,bottom=7pt,
  before skip=9pt,after skip=9pt}
\newtcolorbox{TemplateRecord}[1]{colback=white,colframe=Line,colbacktitle=Pale,
  coltitle=Navy,title={#1},fonttitle=\sffamily\bfseries\small,boxrule=0.5pt,
  arc=1pt,left=9pt,right=9pt,top=6pt,bottom=6pt,before skip=8pt,after skip=8pt}
\pagestyle{fancy}
\fancyhf{}
\fancyhead[L]{\sffamily\scriptsize\color{Muted}\TemplateShortTitle}
\fancyhead[R]{\sffamily\scriptsize\color{Muted}EDITABLE WORKING TEMPLATE}
\fancyfoot[L]{\sffamily\scriptsize\color{Muted}\TemplateCode\enspace\textbullet\enspace EDITION 1.0}
\fancyfoot[R]{\sffamily\scriptsize\color{Muted}\thepage\enspace /\enspace\pageref*{LastPage}}
\renewcommand{\headrulewidth}{0.4pt}
\renewcommand{\footrulewidth}{0pt}
\renewcommand{\headrule}{\hbox to\headwidth{\color{Line}\leaders\hrule height \headrulewidth\hfill}}
\newcommand{\TemplateTitle}[3]{%
  {\sffamily\small\bfseries\color{Teal}ARCHITECTURE DOCUMENTATION\enspace /\enspace #1\par}
  \vspace{14pt}
  {\sffamily\bfseries\fontsize{29}{33}\selectfont\color{Navy}#2\par}
  \vspace{10pt}
  {\sffamily\large\color{Muted}#3\par}
  \vspace{11pt}
  {\sffamily\small\color{Muted}Polished working edition 1.0\enspace\textbullet\enspace 6 September 2026\par}
  \vspace{13pt}{\color{Teal}\hrule height 1.2pt}\vspace{9pt}}
\newcommand{\TemplateAttribution}{%
  \textbf{Attribution.} Adapted from the \emph{Architecture Description Template}
  and \emph{Architecture Viewpoint Template} by Rich Hilliard,
  \copyright\ 2012--2014, supplied under
  \href{https://creativecommons.org/licenses/by/3.0/}{Creative Commons Attribution 3.0 Unported}.
  This adaptation revises layout, wording, fields, dependencies, and edition guidance;
  it does not imply endorsement by the original author or the standards bodies.}
\newcommand{\EditionNote}{%
  The supplied templates were written for ISO/IEC/IEEE 42010:2011.
  This adaptation uses 2022 terminology and adds records for perspectives,
  aspects, and view components. It is an authoring aid, not an official template
  or a complete normative checklist. Assess the completed work against the
  published edition before making a conformance claim.}
\newcommand{\TemplateSources}{%
  {\small
  \begin{enumerate}[label={[\arabic*]},leftmargin=1.6em,itemsep=6pt]
  \item Rich Hilliard. \emph{Architecture Description Template} and
  \emph{Architecture Viewpoint Template}, revision 2.2 (2014), with the supplied
  viewpoint section marked version 2.1b.
  \href{http://www.iso-architecture.org/42010/templates/}{Original source location}.
  The attached source files are the adaptation baseline.
  \item ISO. \emph{ISO/IEC/IEEE 42010:2022: Software, systems and enterprise---Architecture description}.
  \href{https://www.iso.org/standard/74393.html}{Official catalogue and scope}.
  \item ISO/IEC/IEEE. \emph{42010:2022, official preview}.
  \href{https://webstore.ansi.org/preview-pages/iso/preview_iso\%2Biec\%2Bieee\%2B42010-2022.pdf}{Contents and edition changes, hosted by ANSI}.
  \item Creative Commons. \href{https://creativecommons.org/licenses/by/3.0/}{Attribution 3.0 Unported license}.
  \end{enumerate}}}
\renewcommand{\TemplateCode}{AD-TEMPLATE}
\renewcommand{\TemplateShortTitle}{ARCHITECTURE DESCRIPTION TEMPLATE}
\title{Architecture Description\\Template}
\author{Adapted from templates by Rich Hilliard}
\date{6 September 2026}
\hypersetup{pdftitle={Architecture Description Template},pdfauthor={Adapted from templates by Rich Hilliard},pdfsubject={Editable working template with 2022 edition guidance}}
\makeatletter
\renewcommand{\maketitle}{\TemplateTitle{DESCRIPTION}{\@title}{Editable outline and architecture records}}
\makeatother
\begin{document}
\maketitle
\Lead{A reusable outline for describing an architecture, its stakeholder concerns,
its representations, and the evidence behind its decisions.}

\begin{TemplateNote}{Edition and use}
\EditionNote
\end{TemplateNote}

\subsection{Working conventions}
Replace every \FillIn{bracketed field} with project content. Repeat records for
each relevant item. Keep stable identifiers when names or versions change, and
reference controlled artifacts instead of duplicating information unnecessarily.

\textbf{Record fields} are authoring prompts. \textbf{Authoring guidance}
explains how to complete them. Neither label is a quotation of a normative
requirement. Optional extensions are identified in the text.

{\small
\begin{tabularx}{\linewidth}{@{}L{1.7in}Y@{}}
\toprule
\TableHead{Part of the template} & \TableHead{What to establish}\\
\midrule
Identity and context & What architecture is described, for which entity, and at which baseline.\\
Stakeholders and coverage & Whose questions are addressed and where the answers can be found.\\
Viewpoints and views & Conventions for representations and the architecture content that follows them.\\
Consistency and decisions & Relationships, unresolved issues, alternatives, rationale, and review evidence.\\
\bottomrule
\end{tabularx}}

\subsection{Companion template}
Use the \emph{Architecture Viewpoint Template} for a detailed viewpoint
specification. This AD template provides a concise inclusion record and places
for views and view components. A completed AD can reference a controlled
viewpoint specification without repeating its full text.

\vfill
{\footnotesize\color{Muted}\TemplateAttribution\par}

\TemplatePageBreak
\section{Identity, purpose and context}
\Lead{Define the subject and baseline before describing its parts or behavior.}

{\small
\begin{tabularx}{\linewidth}{@{}L{1.52in}Y@{}}
\toprule
\TableHead{Field} & \TableHead{Project entry}\\
\midrule
AD identifier and title & \FillIn{AD-001; descriptive title}\\
Architecture name & \FillIn{Name of the architecture being expressed}\\
Entity of interest & \FillIn{Entity name, type, boundary, and identifying version}\\
Architecture baseline & \FillIn{Current, target, transition, or alternative; effective date}\\
Description status & \FillIn{Draft, reviewed, approved, superseded; revision and issue date}\\
Author and steward & \FillIn{Author, maintaining role, issuing organization, contact route}\\
Review and authority & \FillIn{Reviewers and decision or approval authority, if applicable}\\
Purpose and intended use & \FillIn{Decisions, communication, evaluation, or implementation supported}\\
Configuration reference & \FillIn{Repository, version, tag, controlled document location, or model baseline}\\
Related descriptions & \FillIn{Parent, predecessor, alternative, or companion AD identifiers}\\
\bottomrule
\end{tabularx}}

\subsection{Overview and environment}
\begin{TemplateRecord}{Architecture overview}
\FillIn{Summarize the purpose of the entity, the principal architectural choices,
and the most consequential constraints. Identify the intended audience and
the questions the description will help them answer.}
\end{TemplateRecord}

{\small
\begin{tabularx}{\linewidth}{@{}L{1.52in}Y@{}}
\toprule
\TableHead{Context field} & \TableHead{Project entry}\\
\midrule
Scope and exclusions & \FillIn{Included capabilities, boundaries, lifecycle stages, and explicit exclusions}\\
External entities & \FillIn{People, organizations, services, systems, and physical environment}\\
Interfaces and dependencies & \FillIn{Important exchanges, contracts, dependencies, and ownership boundaries}\\
Assumptions and constraints & \FillIn{Assumptions to validate, constraints, evidence, and responsible roles}\\
\bottomrule
\end{tabularx}}

\subsection{Revision history}
{\small
\begin{tabularx}{\linewidth}{@{}L{0.58in}L{0.9in}Y L{1.1in}@{}}
\toprule
\TableHead{Revision} & \TableHead{Date} & \TableHead{Change and affected baseline} & \TableHead{Author / status}\\
\midrule
\FillIn{1.0} & \FillIn{Date} & \FillIn{Change summary and controlled references} & \FillIn{Role; status}\\
\bottomrule
\end{tabularx}}
\Guide{Add a glossary, reader guide, or contextual diagram where it improves
interpretation. Give the diagram a legend and a source or revision identifier.}

\TemplatePageBreak
\section{Stakeholders and concerns}
\Lead{State concerns as questions the architecture description can help answer.}

\subsection{Stakeholder register}
\Repeat{one row for each individual, group, organization, or stakeholder class.}
{\small
\begin{tabularx}{\linewidth}{@{}L{0.67in}L{1.24in}Y L{1.15in}@{}}
\toprule
\TableHead{ID} & \TableHead{Name / role} & \TableHead{Interest and relationship to the entity} & \TableHead{Concern IDs}\\
\midrule
\RecordID{SH-001} & \FillIn{Stakeholder} & \FillIn{Interest, responsibility, lifecycle involvement} & \FillIn{C-001, ...}\\
\RecordID{SH-002} & \FillIn{Stakeholder} & \FillIn{Interest, responsibility, lifecycle involvement} & \FillIn{C-002, ...}\\
\bottomrule
\end{tabularx}}
\Guide{Consider users, operators, acquirers, owners, suppliers, developers,
builders, and maintainers from the supplied outline. Add other affected parties
appropriate to the entity. Record who was consulted and how their input was captured.}

\subsection{Concern record}
\Repeat{the record below for every concern selected for architectural treatment.}
{\small
\begin{tabularx}{\linewidth}{@{}L{1.49in}Y@{}}
\toprule
\TableHead{Field} & \TableHead{Concern entry}\\
\midrule
Identifier and name & \FillIn{C-001; concise concern name}\\
Question or issue & \FillIn{Specific question, interest, or issue to address}\\
Stakeholders holding it & \FillIn{SH identifiers and relevant differences in expectations}\\
Why it matters & \FillIn{Architectural significance, impact, and lifecycle relevance}\\
Context and priority & \FillIn{Scenario, conditions, priority, source, and date}\\
Assessment basis & \FillIn{Evidence or criterion that would make the answer useful}\\
Coverage references & \FillIn{Relevant perspectives, aspects, viewpoints, and views}\\
Disposition & \FillIn{Addressed, partly addressed, unresolved; explanation and owner}\\
\bottomrule
\end{tabularx}}

\subsection{Prompts for identifying missing concerns}
\begin{itemize}
\item What purposes does the entity serve, and how will architectural suitability
be assessed?
\item What makes construction, acquisition, deployment, and operation feasible?
\item Which risks, external effects, and trade-offs matter to affected stakeholders?
\item How will the entity be maintained, changed, and eventually retired?
\end{itemize}

\begin{TemplateNote}{Evidence is different from a concern statement}
\small
\emph{How does the service recover after a dependency fails?} is a concern.
A recovery scenario, stated assumptions, measured outcome, and assessment
provide evidence addressing it. A named view alone does not establish the answer.
\end{TemplateNote}

\TemplatePageBreak
\section{Perspectives, aspects and traceability}
\Lead{Organize the concerns and make their treatment navigable. These tables are
project record formats, not a prescribed taxonomy.}

\subsection{Stakeholder perspective register}
{\small
\begin{tabularx}{\linewidth}{@{}L{0.52in}L{1.28in}Y L{1.17in}@{}}
\toprule
\TableHead{ID} & \TableHead{Perspective} & \TableHead{Shared orientation and stakeholders} & \TableHead{Concern IDs}\\
\midrule
\RecordID{P-001} & \FillIn{Name} & \FillIn{How concerns are grouped; SH identifiers} & \FillIn{C identifiers}\\
\RecordID{P-002} & \FillIn{Name} & \FillIn{How concerns are grouped; SH identifiers} & \FillIn{C identifiers}\\
\bottomrule
\end{tabularx}}
\Guide{Explain the grouping so readers can distinguish a stakeholder perspective
from a viewpoint specification. Name the concerns that make the grouping useful.}

\subsection{Aspect register}
{\small
\begin{tabularx}{\linewidth}{@{}L{0.52in}L{1.28in}Y L{1.17in}@{}}
\toprule
\TableHead{ID} & \TableHead{Aspect} & \TableHead{Feature of the entity being considered} & \TableHead{Related IDs}\\
\midrule
\RecordID{A-001} & \FillIn{Name} & \FillIn{Scope and relevance to the architecture} & \FillIn{C, VP, V IDs}\\
\RecordID{A-002} & \FillIn{Name} & \FillIn{Scope and relevance to the architecture} & \FillIn{C, VP, V IDs}\\
\bottomrule
\end{tabularx}}
\Guide{Possible organizing choices include structure, behavior, information,
deployment, or evolution. Define the chosen meaning instead of relying on a label.}

\subsection{Concern-to-view coverage matrix}
{\small
\begin{tabularx}{\linewidth}{@{}L{0.64in}L{0.91in}L{0.89in}L{0.86in}Y@{}}
\toprule
\TableHead{Concern} & \TableHead{Stakeholders} & \TableHead{Viewpoints} & \TableHead{Views} & \TableHead{Answer / evidence location}\\
\midrule
\FillIn{C-001} & \FillIn{SH IDs} & \FillIn{VP IDs} & \FillIn{V IDs} & \FillIn{Artifact, section, component, and status}\\
\FillIn{C-002} & \FillIn{SH IDs} & \FillIn{VP IDs} & \FillIn{V IDs} & \FillIn{Artifact, section, component, and status}\\
\bottomrule
\end{tabularx}}

\subsection{Identifier relationships to preserve}
{\small
\begin{tabularx}{\linewidth}{@{}L{1.35in}Y@{}}
\toprule
\TableHead{Relationship} & \TableHead{Meaning in this template}\\
\midrule
SH to C & Which stakeholder holds the concern.\\
VP to C & Which concern the viewpoint is intended to frame.\\
V to VP & Which specification governs the view.\\
VC to V & Which view or views include the component.\\
VC to MK & Which model kind governs a model-based component.\\
CR to AD elements & Which identified items participate in a recorded correspondence.\\
\bottomrule
\end{tabularx}}

\Guide{Review missing links, ambiguous references, and incomplete concern
treatment. Keep unresolved items visible and assigned rather than marking
coverage complete because a diagram or section exists.}

\TemplatePageBreak
\section{Viewpoint selection and specifications}
\Lead{Select viewpoints for the concerns they frame and the work their views support.}

\subsection{Viewpoint catalogue}
{\small
\begin{tabularx}{\linewidth}{@{}L{0.67in}L{1.45in}Y L{1.05in}@{}}
\toprule
\TableHead{ID} & \TableHead{Name and version} & \TableHead{Concern coverage and selection rationale} & \TableHead{Governs}\\
\midrule
\RecordID{VP-001} & \FillIn{Name; version} & \FillIn{C IDs; why this viewpoint is useful} & \FillIn{V IDs}\\
\RecordID{VP-002} & \FillIn{Name; version} & \FillIn{C IDs; why this viewpoint is useful} & \FillIn{V IDs}\\
\bottomrule
\end{tabularx}}

\subsection{Viewpoint inclusion record}
\Repeat{for each selected viewpoint; include its specification or a controlled reference.}
{\small
\begin{tabularx}{\linewidth}{@{}L{1.47in}Y@{}}
\toprule
\TableHead{Field} & \TableHead{Viewpoint entry}\\
\midrule
Identity and source & \FillIn{VP identifier, name, version, source, owner, and locator}\\
Purpose and applicability & \FillIn{Intended use, relevant situations, and limitations}\\
Audience and coverage & \FillIn{Typical stakeholders, concerns, perspectives, and aspects}\\
Model kinds & \FillIn{MK identifiers; conventions, notations, and versions}\\
Other representations & \FillIn{Non-model-based components and interpretation legends}\\
View methods & \FillIn{Construction, interpretation, analysis, and other defined methods}\\
Correspondence methods & \FillIn{Relationships to establish and methods for checking them}\\
Selection rationale & \FillIn{Why selected, alternatives considered, and expected contribution}\\
Specification location & \FillIn{Embedded specification or exact external artifact and version}\\
\bottomrule
\end{tabularx}}

\subsection{Conventions worth making explicit}
\begin{itemize}
\item Define what each symbol, relationship, table column, and direction means.
State the scope, abstraction level, and assumptions of each representation.
\item Separate rules about the modeling language from constraints on the entity
being described. Give both a source and an interpretation.
\item State how models or other view components will be built, read, and assessed.
Identify checks that depend on another view or an external information source.
\end{itemize}

\begin{TemplateNote}{Use the companion viewpoint document}
\small
The companion guide and worksheet provide detailed records for concerns,
model kinds, view methods, correspondences, examples, and sources. The viewpoint
defines conventions. Its views contain architecture-specific information
prepared using those conventions.
\end{TemplateNote}

\TemplatePageBreak
\section{Views and view components}
\Lead{Repeat the view record for each view, then record its included components.}

\subsection{View record}
{\small
\begin{tabularx}{\linewidth}{@{}L{1.41in}Y@{}}
\toprule
\TableHead{Field} & \TableHead{Architecture view entry}\\
\midrule
Identity and revision & \FillIn{V identifier, title, owner, revision, and status}\\
Governing viewpoint & \FillIn{VP identifier and exact specification version}\\
Scope and baseline & \FillIn{Entity coverage, architecture baseline, assumptions, exclusions}\\
Coverage and purpose & \FillIn{C, P, and A identifiers; decisions or questions supported}\\
Included components & \FillIn{VC identifiers and locations, including shared components}\\
Interpretation and findings & \FillIn{Principal observations, conclusions, and relevant evidence}\\
Known issues & \FillIn{Missing content, convention deviations, contradictions, and issue IDs}\\
\bottomrule
\end{tabularx}}

\subsection{View-component record}
\Repeat{for every diagram, model-derived representation, table, narrative, or other component.}
{\small
\begin{tabularx}{\linewidth}{@{}L{1.41in}Y@{}}
\toprule
\TableHead{Field} & \TableHead{View-component entry}\\
\midrule
Identity and ownership & \FillIn{VC identifier, name, author or owner, and revision}\\
Membership & \FillIn{V identifiers for all views that use this component}\\
Source and locator & \FillIn{Source artifact or model, version, location, and derivation}\\
Representation type & \FillIn{Model-based or non-model-based; descriptive format}\\
Governing model kind & \FillIn{MK identifier and version, when model-based}\\
Legend and conventions & \FillIn{Symbols, relationships, labels, units, direction, and interpretation}\\
Architectural content & \FillIn{Insert the actual representation or a controlled reference}\\
Assumptions and limitations & \FillIn{Scope, data quality, abstraction, uncertainty, and known deviations}\\
\bottomrule
\end{tabularx}}

\Guide{The 2022 edition accommodates both model-based and non-model-based
view components. Provide a legend for interpretation; identify the model kind
for model-based content. See the edition preview in the references.}

\subsection{Concern assessment within the view}
{\small
\begin{tabularx}{\linewidth}{@{}L{0.64in}Y L{1.18in}@{}}
\toprule
\TableHead{Concern} & \TableHead{Answer, evidence, and relevant VC identifiers} & \TableHead{Limit / next step}\\
\midrule
\FillIn{C-001} & \FillIn{What the view establishes and the evidence supporting it} & \FillIn{Gap; owner; action}\\
\bottomrule
\end{tabularx}}

\TemplatePageBreak
\section{Correspondences and consistency}
\Lead{Record relationships separately from the methods used to establish or assess them.}

\subsection{Correspondence register}
{\small
\begin{tabularx}{\linewidth}{@{}L{0.62in}L{1.48in}Y L{0.95in}@{}}
\toprule
\TableHead{ID} & \TableHead{Participating elements} & \TableHead{Relation and meaning} & \TableHead{Method IDs}\\
\midrule
\RecordID{CR-001} & \FillIn{Two or more IDs; include AD IDs if external} & \FillIn{Identity, allocation, refinement, dependency, or other named relation} & \FillIn{CM IDs}\\
\RecordID{CR-002} & \FillIn{Participating IDs} & \FillIn{Meaning and rationale} & \FillIn{CM IDs}\\
\bottomrule
\end{tabularx}}

\subsection{Correspondence-method record}
\Repeat{for each method used to establish, check, or maintain relationships.}
{\small
\begin{tabularx}{\linewidth}{@{}L{1.35in}Y@{}}
\toprule
\TableHead{Field} & \TableHead{Method entry}\\
\midrule
Identity and source & \FillIn{CM identifier, name, version, origin, and owner}\\
Scope and inputs & \FillIn{Applicable relations, AD elements, baselines, and input sources}\\
Rule or procedure & \FillIn{Relationship to establish; checking steps; interpretation of outcomes}\\
Execution and evidence & \FillIn{Reviewer or tool, execution date, evidence location, and baseline}\\
Result and violations & \FillIn{Satisfied, violated, inconclusive, or not assessed; affected IDs}\\
\bottomrule
\end{tabularx}}

\begin{TemplateNote}{Illustrative distinction}
\small
\RecordID{CR-001} can record that two view elements represent the same service.
\RecordID{CM-001} can check that their identifiers, responsibilities, and baseline
agree. Recording the relationship does not, by itself, demonstrate that the
check has run or succeeded.
\end{TemplateNote}

\subsection{Inconsistency and unresolved-issue record}
{\small
\begin{tabularx}{\linewidth}{@{}L{1.35in}Y@{}}
\toprule
\TableHead{Field} & \TableHead{Issue entry}\\
\midrule
Identity and type & \FillIn{ISS identifier; contradiction, missing information, deviation, or open question}\\
Affected elements & \FillIn{AD, VP, V, VC, MK, C, or other relevant identifiers}\\
Description and impact & \FillIn{What conflicts or remains unknown; consequences for interpretation}\\
Disposition and owner & \FillIn{Proposed action, responsible role, decision needed, and target date}\\
Closure evidence & \FillIn{Resolution, supporting record, date, and review status}\\
\bottomrule
\end{tabularx}}
\Guide{Distinguish a recorded inconsistency from an unassessed relationship.
Use explicit status values so readers do not mistake an empty field for a successful check.}

\TemplatePageBreak
\section{Architecture decisions, rationale and evaluation}
\Lead{Preserve the reasons for a choice so future changes can be assessed against its context.}

\subsection{Architecture decision record}
\Repeat{for each decision selected for inclusion in this AD.}
{\small
\begin{tabularx}{\linewidth}{@{}L{1.39in}Y@{}}
\toprule
\TableHead{Field} & \TableHead{Decision entry}\\
\midrule
Identity and status & \FillIn{ADR identifier, title, date, proposed / accepted / superseded}\\
Owner and authority & \FillIn{Decision owner, participants, and authority where applicable}\\
Context and forces & \FillIn{Problem, constraints, risks, assumptions, and relevant baseline}\\
Concern links & \FillIn{C identifiers and stakeholder interests affected}\\
Decision statement & \FillIn{What is selected and the scope of that choice}\\
Rationale & \FillIn{Reasoning and evidence that support the selected alternative}\\
Affected elements & \FillIn{AD element identifiers and external dependencies}\\
Consequences & \FillIn{Benefits, costs, disadvantages, obligations, and residual uncertainty}\\
Revisit conditions & \FillIn{Events or changed assumptions that warrant reassessment}\\
\bottomrule
\end{tabularx}}

\subsection{Alternatives and trade-offs}
{\small
\begin{tabularx}{\linewidth}{@{}L{1.0in}Y Y L{0.93in}@{}}
\toprule
\TableHead{Alternative} & \TableHead{Benefits and supporting evidence} & \TableHead{Costs, constraints, and uncertainty} & \TableHead{Disposition}\\
\midrule
\FillIn{Option A} & \FillIn{Evidence and assumptions} & \FillIn{Trade-offs and limits} & \FillIn{Selected / other}\\
\FillIn{Option B} & \FillIn{Evidence and assumptions} & \FillIn{Trade-offs and limits} & \FillIn{Reason}\\
\bottomrule
\end{tabularx}}

\subsection{Architecture evaluation record}
{\small
\begin{tabularx}{\linewidth}{@{}L{1.39in}Y@{}}
\toprule
\TableHead{Field} & \TableHead{Evaluation entry}\\
\midrule
Evaluation identity & \FillIn{EV identifier, date, participants, and evaluated baseline}\\
Questions and method & \FillIn{Concern IDs, criteria, scenarios, inputs, and procedure}\\
Results and confidence & \FillIn{Findings, evidence, assumptions, and limitations}\\
Actions and traceability & \FillIn{ADR or ISS links, responsible roles, and follow-up actions}\\
\bottomrule
\end{tabularx}}
\Guide{Use measured results where available. Identify estimates and hypotheses
as such, and avoid presenting an unevaluated architectural preference as an established benefit.}

\TemplatePageBreak
\section{Review, references and provenance}
\Lead{Use the records below to prepare the completed description for review.}

\subsection{Review navigation}
{\small
\begin{tabularx}{\linewidth}{@{}L{0.82in}L{1.1in}Y@{}}
\toprule
\TableHead{2022 area} & \TableHead{Template part} & \TableHead{Review focus}\\
\midrule
6.1 & Section 1 & Identity, context, purpose, and description baseline.\\
6.2--6.5 & Sections 2--3 & Stakeholders, perspectives, concerns, aspects, and coverage.\\
6.6--6.8 & Sections 4--5 & Viewpoints, views, view components, and interpretation.\\
6.9 & Section 6 & Correspondences, their methods, and known inconsistencies.\\
6.10 & Section 7 & Decisions, reasoning, alternatives, and supporting evidence.\\
\bottomrule
\end{tabularx}}
{\footnotesize\color{Muted}Clause areas are navigation references from the published
contents, not a complete transcription or assessment of the requirements.\par}

\subsection{Requirement-to-evidence record}
{\small
\begin{tabularx}{\linewidth}{@{}L{1.04in}Y L{1.17in}@{}}
\toprule
\TableHead{Requirement ref.} & \TableHead{AD location, evidence, baseline, and reviewer} & \TableHead{Result / gap}\\
\midrule
\FillIn{Edition; clause} & \FillIn{Precise location and supporting artifact} & \FillIn{Result; ISS link}\\
\bottomrule
\end{tabularx}}

\textbf{Review prompts.} Confirm resolved identifiers, accessible references,
consistent baselines, interpretable legends, concern treatment, decision rationale,
and explicit issue status. Replace template fields before issuing a project AD;
keep the underlying evidence available to its intended readers.

\subsection{References and adaptation history}
\TemplateSources
{\footnotesize\color{Muted}\TemplateAttribution\par
\textbf{Changes in this edition:} new typography and records; self-contained build;
removal of missing image and bibliography dependencies; revised 2022 terminology;
additional coverage and evidence fields. The supplied 2011-era normative labels
have been replaced with clearly identified authoring guidance.\par}

\end{document}
